\section{Verifiche numeriche complementari}
\label{app:controls}
\setcounter{table}{0}\setcounter{equation}{0}

\subsection{Convergenza, energia e campionamento}
La Tabella~\ref{tab:mms-compact} riporta gli errori in norma $L^2$ delle due verifiche con soluzioni manufatte. Il raffinamento congiunto riduce il passo spaziale e quello temporale. I tassi indicati nei Risultati descrivono le sequenze provate.

\begin{table}[!htbp]
\centering\small
\caption{Errori $L^2$ al tempo finale. Le prime tre righe usano elementi $[Q_2]^2/Q_1/Q_1$, le altre elementi quadratici per tutti i campi.}
\label{tab:mms-compact}
\begin{tabular}{@{}rrccc@{}}
\toprule
$N$ & $M$ & Spostamento & Pressione & Concentrazione\\
\midrule
4  & 100  & $2.6868\times10^{-3}$ & $5.4526\times10^{-3}$ & $2.4562\times10^{-2}$\\
8  & 400  & $3.5294\times10^{-4}$ & $1.3055\times10^{-3}$ & $6.3959\times10^{-3}$\\
16 & 1600 & $4.4932\times10^{-5}$ & $3.2286\times10^{-4}$ & $1.6164\times10^{-3}$\\
\addlinespace
24 & 10   & $1.3569\times10^{-3}$ & $2.7205\times10^{-2}$ & $1.3713\times10^{-1}$\\
24 & 20   & $6.3341\times10^{-4}$ & $1.3843\times10^{-2}$ & $7.0004\times10^{-2}$\\
24 & 40   & $3.0010\times10^{-4}$ & $6.9863\times10^{-3}$ & $3.5421\times10^{-2}$\\
24 & 80   & $1.4436\times10^{-4}$ & $3.5106\times10^{-3}$ & $1.7835\times10^{-2}$\\
\bottomrule
\end{tabular}
\end{table}

Sia $R_n$ il residuo cumulativo del bilancio energetico discreto, ottenuto sommando variazione di energia e dissipazioni e sottraendo il lavoro esterno. Testando l'equilibrio con $\Delta U^n$ e il bilancio del fluido con $\Delta t_n P^n$, si ottiene
\[
 R_n-R_{n-1}
 =-\tfrac12(\Delta U^n)^{\mathsf T}K_E\Delta U^n
  -\tfrac12S_p(\Delta P^n)^{\mathsf T}M_p\Delta P^n .
\]
Lo scarto massimo di questa identità è $4.78\times10^{-17}$. Il residuo cumulativo in norma riflette quindi anche la dissipazione numerica del passo. Per il controllo della massa del tracciante si considera la quantità pesata con $\phi$, in linea con l'uso di bilanci come diagnostica per equazioni frazionarie~\citep{cardone2022numerical}.

Sul quadrato, la valutazione $Q_1$ e l'interpolazione su triangolazione differiscono dello $0.0360\%$ sul target e dello $0.1332\%$ sulla griglia inversa. La sostituzione varia il Jacobiano dell'$1.031\%$ in norma di Frobenius e i singoli valori singolari al più dello $0.136\%$. In anatomia, la mesh a 3~mm ha Dice $0.9982$ rispetto alla maschera esterna. Il Jacobiano scalato minimo delle celle vale $0.083$. Dimezzando il passo delle differenze finite da $10^{-3}$ a $5\times10^{-4}$, la variazione relativa di $J_6$ vale $1.24\times10^{-6}$ nel disegno bilanciato e $9.10\times10^{-7}$ nell'edema.

\subsection{Recupero cartesiano e peso del prior}
Sul quadrato, con $s_{\rm ref}=\|\mathbf y^\dagger\|_\infty$, i dati puliti usano $W_A=I/(\sqrt{N_{\rm obs}}\,s_{\rm ref})$ e $\lambda_A=10^{-4}$ con centro $\mathbf z_0=(0.5,0.5,0.5,0.5)^{\mathsf T}$. Per le osservazioni rumorose usiamo $W_B=I/\sigma$ con $\sigma=\eta_\infty s_{\rm ref}$. Il termine dati è quindi moltiplicato per $N_{\rm obs}\eta_\infty^{-2}$ rispetto a $W_A$. A parità di residuo, una $\sigma$ doppia richiede una $\lambda_B$ quattro volte minore per conservare il peso relativo della penalizzazione.

Usiamo il Particle Swarm Optimization continuo (CPSO)~\citep{kennedyndparticle,orlando2021analyticCPSO} e lo confrontiamo con Differential Evolution e Powell sullo stesso target pulito, con al più 600 valutazioni del problema diretto per esecuzione~\citep{storn1997differential,powell1964an}. I migliori costi sono simili, mentre $\widehat\beta$ con CPSO varia da $0.3880$ a $0.6033$ su tre partenze. Nella scansione di profilo su 15 nodi il minimo campionato passa da $0.4500$ per il solo termine dati a $0.479143$ includendo la penalizzazione sulla griglia inversa $(16,80)$. Per il solo termine dati, sulla griglia $(24,120)$ il minimo campionato è a $0.4375$. I limiti di valutazione raggiunti in diverse riottimizzazioni rendono questi risultati diagnostiche del protocollo.
